\documentclass[11pt]{article}
\usepackage[margin=21mm]{geometry}
\usepackage{lmodern}
\usepackage[T1]{fontenc}
\usepackage{amsmath}
\usepackage[strings]{underscore}
\usepackage{textcomp}
\usepackage{microtype}
\usepackage{xcolor}
\usepackage{booktabs}
\usepackage{enumitem}
\usepackage{titlesec}
\usepackage[colorlinks=true,linkcolor=tealc,urlcolor=tealc]{hyperref}

\definecolor{coral}{HTML}{B85B3A}
\definecolor{tealc}{HTML}{2A6478}
\definecolor{forest}{HTML}{3D6E4A}
\definecolor{char}{HTML}{2A2520}
\definecolor{dimc}{HTML}{6E6862}
\definecolor{sand}{HTML}{F5F1EA}
\definecolor{tealL}{HTML}{D6E3E8}

\setlist[itemize]{leftmargin=1.3em, itemsep=2pt, topsep=2pt}
\titlespacing*{\section}{0pt}{12pt}{5pt}
\titlespacing*{\subsection}{0pt}{9pt}{3pt}
\titleformat{\section}{\Large\bfseries\color{char}}{\thesection}{0.6em}{}
\titleformat{\subsection}{\normalsize\bfseries\color{tealc}}{}{0em}{}
\newcommand{\code}[1]{\texttt{\small #1}}
\newcommand{\fn}[1]{\texttt{\small\color{coral}#1}}
\newcommand{\arr}{\textrightarrow}
\newcommand{\nb}[1]{\texttt{\small\color{forest}#1}}

\begin{document}
\sloppy

\begin{center}
  {\LARGE\bfseries\color{char} Notebook guide}\\[3pt]
  {\color{dimc} How to run everything, and what code each part uses.}
\end{center}
\vspace{2pt}

\noindent
\textbf{Keep this open while you work.} It maps every notebook --- and each step
\emph{inside} a notebook --- to the exact \code{.py} file behind it. Two
references live in the \code{code/notebooks/} folder:

\begin{center}\small
\begin{tabular}{@{}ll@{}}
\toprule
File & Use it for \\
\midrule
\textbf{NOTEBOOK\_GUIDE} (this) & how to \emph{run} the notebooks + which code each step uses \\
\textbf{CODE\_MANUAL} & the \emph{file-by-file} reference: what every \code{.py} does + how they connect \\
\bottomrule
\end{tabular}
\end{center}

\noindent
Rule of thumb: \textbf{this guide tells you \emph{which} code a cell runs; the
code manual tells you \emph{what that code does inside}.}

% =============================================================================
\section{Running them in Jupyter}

\begin{quote}\ttfamily\small
cd Lunar-HFE\\
python3 -m venv .venv \&\& source .venv/bin/activate\\
pip install -e "code[dev]"\hfill\textrm{\itshape\# installs the \code{lunar} package}\\
jupyter lab code/notebooks/
\end{quote}

Open the notebooks \textbf{in order, 00 \arr\ 06}, and use \emph{Run All} on
each. Every notebook finds the repo by itself
(\fn{lunar.\_bootstrap.find\_repo\_root()}), so it works no matter where Jupyter
was launched. The results are \textbf{cached} --- you can read the numbers and
figures without waiting for the long solves to finish.

% =============================================================================
\section{The three kinds of code the notebooks call}

Every cell touches one of three layers. Knowing which tells you what it does:

\begin{center}\footnotesize
\begin{tabular}{@{}lll@{}}
\toprule
Layer & Folder & What it is \\
\midrule
\textbf{Engine} & \code{code/src/lunar/} & the physics \emph{library} (solver, equilibrium, grid, config). Everything imports it. \\
\textbf{Pipeline} & \code{code/pipeline/} & the \emph{application}: \code{compute/} \arr\ \code{results/*.json}; \code{figures/} \arr\ PDFs. Writes files. \\
\textbf{Tests} & \code{code/tests/} & pytest checks the engine is correct. \\
\bottomrule
\end{tabular}
\end{center}

\noindent
So \code{import lunar.equilibrium} calls the \textbf{engine}; \code{import
make\_letter\_figures} or running \code{retrieve\_kd.py} calls the
\textbf{pipeline}.

% =============================================================================
\section{Reading order at a glance}

The order is \textbf{chronological}: set up, see what we measure, learn the
solver, run it, then the result and discussion figures, and finally performance.

\begin{center}\footnotesize
\begin{tabular}{@{}clll@{}}
\toprule
\# & Notebook & Produces & Main code it uses \\
\midrule
00 & \nb{00\_setup}      & environment + data-load check & \code{lunar.apollo\_helpers, ephem} \\
01 & \nb{01\_methods}    & letter Figs 1--3, 5 & \code{pipeline/figures/} (4 generators) \\
02 & \nb{02\_anchor\_method}     & flux-anchored solver, explained & \code{lunar.equilibrium}, \code{make\_equilibrium\_demo} \\
03 & \nb{03\_retrieval}  & \textbf{core K\_d retrieval} \arr\ JSON & \code{pipeline/compute/retrieve\_kd.py} \\
04 & \nb{04\_results}    & letter Figs 6--10 & \code{make\_results\_figures}, \code{compute\_diviner\_closure} \\
05 & \nb{05\_discussion} & letter Figs 11--12 & \code{make\_prior\_estimates\_figure} \\
06 & \nb{06\_performance}& C++ kernel: agree, speed, live retrieval & \code{benchmark\_cpp}, \code{retrieve\_kd} \\
\bottomrule
\end{tabular}
\end{center}

\noindent
\textbf{One linear read:} each notebook builds on the previous --- you meet the
solver (\nb{02}) before the retrieval that uses it (\nb{03}), and see the result
before the discussion. The two \emph{explainers} are \nb{02} (how the solver
works) and \nb{06} (how the backends compare). The one hard rule: run
\nb{03\_retrieval} before \nb{04}/\nb{05}, which read its JSON.

% =============================================================================
\section{Inside each notebook: step \arr\ code \arr\ what happens}

\subsection{\nb{00\_setup} --- is everything wired up?}
\begin{itemize}
\item environment check \arr\ \fn{lunar.\_bootstrap.find\_repo\_root()} locates \code{code/} and confirms \code{lunar} imports.
\item load the Apollo record \arr\ \fn{lunar.apollo\_helpers.extract\_sensor\_stability()} (reads \code{code/data/apollo/} via \code{lunar.validation}).
\item solar geometry \arr\ \code{lunar.ephem} (SPICE kernels).
\item \emph{(one-time)} \fn{pipeline/fetch\_diviner.py} downloads the Diviner tiles used by Fig 10.
\end{itemize}

\subsection{\nb{01\_methods} --- the methods figures (Figs 1--3, 5)}
Each cell: \textbf{import a generator \arr\ call it (writes \code{figures/<name>.pdf}) \arr\ \code{show\_fig()} displays it.}
\begin{itemize}
\item Fig 1 probe schematic \arr\ \fn{make\_intro\_figures.fig\_intro\_probe()} --- both borestems from the real sensor depths (\code{lunar.apollo\_helpers}) + the K(T,z) comparison (\code{lunar.properties}).
\item Fig 2 context map \arr\ \fn{make\_context\_map\_figure.main()} --- three nearside globes; uses the LOLA DEM in \code{code/data/lola/}.
\item Fig 3 sensor timeline \arr\ \fn{make\_apollo\_timeline\_letter.main()} --- the per-sensor stability windows.
\item Fig 5 amplitude vs depth \arr\ \fn{make\_letter\_figures.fig\_amplitude\_vs\_depth()} --- diurnal-wave decay (\code{lunar.solver}).
\end{itemize}

\subsection{\nb{02\_anchor\_method} --- how the flux-anchored solver works}
The heart of the study, explained before you use it in \nb{03}.
\begin{itemize}
\item builds the depth grid \arr\ \fn{lunar.grid.make\_geometric\_grid}
\item brute-force spin-up vs the fast flux-anchored solve, shown to reach the \emph{same} steady state \arr\ \fn{lunar.equilibrium.solve\_periodic\_equilibrium}, \code{lunar.solver}
\item reading temperatures off a converged cycle \arr\ \fn{lunar.equilibrium.profile\_at\_time} / \code{profile\_at\_local\_time}
\item the teaching animations \arr\ \fn{make\_equilibrium\_demo}
\end{itemize}

\subsection{\nb{03\_retrieval} --- the core computation (writes the headline JSON)}
\begin{itemize}
\item \textbf{Step 1} runs \fn{pipeline/compute/retrieve\_kd.py} (subprocess). Inside, per site it: loads the fit targets (\code{lunar.apollo\_helpers}), \textbf{sweeps K\_d} solving the model to steady state each time with the method from \nb{02} (\code{lunar.equilibrium.solve\_periodic\_equilibrium} \arr\ \code{lunar.solver}), takes the RMSE minimum (\textbf{K\_d*}), and \textbf{bootstraps} the CI. Writes \code{code/results/kd\_retrieval\_results.json} --- \emph{the file every figure notebook reads.}
\item \textbf{verify cell} \arr\ reads that JSON and prints K\_d* and the 95\% CI.
\item \textbf{Step 2 (optional)} \arr\ the \code{compute\_*} scripts; each re-runs the retrieval under one perturbed assumption (Table 3 error budget).
\item \textbf{Step 3 (optional)} \arr\ \fn{bayesian\_crosscheck.py} --- an independent \code{emcee} MCMC cross-check.
\end{itemize}

\subsection{\nb{04\_results} --- the results figures (Figs 6--10)}
Same generator \arr\ \code{show\_fig()} pattern; these read the JSON from \nb{03}:
\begin{itemize}
\item Fig 6 mean-T profile \arr\ \fn{make\_letter\_figures.fig\_mean\_T\_profile()}
\item Fig 7 K\_d sweep \arr\ \fn{make\_letter\_figures.fig\_kd\_sweep()}
\item Fig 8 bootstrap \arr\ \fn{make\_results\_figures.fig\_bootstrap()}
\item Fig 9 thermal profiles \arr\ \fn{make\_results\_figures.fig\_thermal\_profiles()}
\item Fig 10 Diviner closure \arr\ runs \fn{pipeline/compute/compute\_diviner\_closure.py} (model vs Diviner; needs the fetched data).
\end{itemize}

\subsection{\nb{05\_discussion} --- the discussion figures (Figs 11--12)}
\begin{itemize}
\item Fig 11 robustness \arr\ \fn{make\_results\_figures.fig\_robustness()} --- the inter-site contrast under stress tests.
\item Fig 12 five-decade K\_d context \arr\ \fn{make\_prior\_estimates\_figure.main()} --- the two-panel figure separating contact-asymptote K\_d from effective K.
\end{itemize}

\subsection{\nb{06\_performance} --- the C++ kernel and the retrieval it powers}
All about the one demanding primitive (the solver march):
\begin{itemize}
\item \textbf{\S1 agreement} compiles \code{code/cpp/solver.cpp} (\code{clang++}) \emph{in the notebook} and runs \emph{all three} engines --- generic Python, Numba, C++ --- on identical inputs through \fn{pipeline/compute/benchmark\_cpp.py} (\code{dump\_inputs / run\_cpp / run\_py}), confirming they agree to $\sim$10$^{-12}$ K.
\item \textbf{\S2 speed} plots the steady-state ms/lunation per engine from \code{results/cpp\_benchmark.json}.
\item \textbf{\S3 the real workload} runs the actual \fn{retrieve\_kd.run\_kd\_sweep\_extended} \emph{live with a \code{tqdm} progress bar}, lands on the published K\_d* (4.600 / 7.079), and shows the backend cost of the full 61-solve sweep.
\end{itemize}

\begin{center}\small\itshape\color{dimc}
For what any of these modules do internally --- every function and how they
depend on one another --- see \textbf{CODE\_MANUAL.pdf} in this folder.
\end{center}

% =============================================================================
\section{The other running order: \texttt{make all}}

The notebooks \emph{narrate} the science; the \textbf{Makefile} regenerates
everything from scratch:

\begin{enumerate}[leftmargin=1.6em]
\item \code{make retrieve} \arr\ \fn{retrieve\_kd.py} \arr\ \code{kd\_retrieval\_results.json}
\item \code{make aux} \arr\ the sensitivity / model-selection / error-budget sweeps \arr\ more \code{results/*.json}
\item \code{make figures} \arr\ \fn{pipeline/make\_all\_figures.py} runs every generator \arr\ \code{figures/}
\item \code{make paper} \arr\ compiles the documents in \code{documents/}
\end{enumerate}

\end{document}
